% GENERATED FILE. Edit canonical lessons, then run npm run generate:books.
% canonical-source-hash: fnv1a64:3ed130c56b84e28e
% canonical-lessons: TA-C469-pakai, TA-C469-virotam, TA-C469-vancam, TA-C469-etiri, TA-C469-palivanku

\chapter{Enemies and Grudges}
\label{ch:ta-enemies-and-grudges}
\hlchaptermodality{469}

\begin{chapteropening}
\textbf{By the end of this chapter:} \emph{I can talk about enmity, grudges and revenge.}

Complete the last lesson of chapter 469: I can talk about enmity, grudges and revenge.
\end{chapteropening}

\section[\texorpdfstring{\ta{பகை} (pakai)}{பகை (pakai)}]{\ta{பகை} (pakai) — enmity; a feud}

\label{lesson:TA-C469-pakai}



\begin{quote}
Before the new one: say the Tamil for to weaken, then the Tamil for to nod off.
\end{quote}

\subsection*{You'll want to know: \ta{பகை}}
\textbf{\ta{பகை}} — \emph{pakai} — \textquotedblleft{}enmity; a feud\textquotedblright{}.

\begin{quote}
\ta{இரண்டு} \ta{குடும்பங்களுக்கும்} \ta{இடையில்} \ta{பல} \ta{வருடமாகப்} \ta{பகை}.
\end{quote}

\begin{quote}
\emph{iraṇṭu kuṭumpaṅkaḷukkum iṭaiyil pala varuṭamākap pakai.}
\end{quote}

\begin{quote}
There has been a feud between the two families for many years.
\end{quote}

\ta{பகைவன்}, an enemy, in written Tamil.

\begin{cousinweb}
\textbf{\ta{பகை}} (\emph{pakai}) is native Dravidian (DEDR 3808), in the family of \ta{பகு} \emph{paku}, \textquotedblleft{}to be split, to be at variance\textquotedblright{}.
\end{cousinweb}

\begin{grammarlens}[title={What you've built}]
One more word for describing what you lived and felt.
\end{grammarlens}

\subsection*{Your turn}
\begin{itemize}
\raggedright
  \item \emph{Say it:} \emph{pakai}
  \item \emph{Say it:} \emph{pakai}, once more
  \item \emph{Say it:} \emph{iraṇṭu kuṭumpaṅkaḷukkum iṭaiyil pala varuṭamākap pakai.} — There has been a feud between the two families for many years.
  \item \emph{Recall:} say the Tamil for to weaken, then the Tamil for to nod off, then say \emph{pakai} again
\end{itemize}

\begin{tcolorbox}[breakable,colback=teal!4,colframe=teal!35!black,title={Before you move on}]
What does \ta{பகை} mean? (\textquotedblleft{}Enmity; a feud\textquotedblright{}.) Say it once more.
\end{tcolorbox}
\section[\texorpdfstring{\ta{விரோதம்} (virōtam)}{விரோதம் (virōtam)}]{\ta{விரோதம்} (virōtam) — hostility; ill will}

\label{lesson:TA-C469-virotam}



\begin{quote}
Before the new one: say the Tamil for to nod off, then the Tamil for enmity.
\end{quote}

\subsection*{You'll want to know: \ta{விரோதம்}}
\textbf{\ta{விரோதம்}} — \emph{virōtam} — \textquotedblleft{}hostility; ill will\textquotedblright{}.

\begin{quote}
\ta{ஒரு} \ta{சிறிய} \ta{சண்டையால்} \ta{அவர்களுக்குள்} \ta{விரோதம்} \ta{வளர்ந்தது}.
\end{quote}

\begin{quote}
\emph{oru ciṟiya caṇṭaiyāl avarkaḷukkuḷ virōtam vaḷarntatu.}
\end{quote}

\begin{quote}
Because of a small quarrel, hostility grew between them.
\end{quote}

-\ta{உக்குள்} \textquotedblleft{}within, between\textquotedblright{}: \ta{அவர்களுக்குள்}, between them.

\begin{cousinweb}
\textbf{\ta{விரோதம்}} (\emph{virōtam}) is borrowed from Sanskrit \emph{virōdha}, \textquotedblleft{}opposition, hostility\textquotedblright{} (Monier-Williams).
\end{cousinweb}

\begin{grammarlens}[title={What you've built}]
One more word for describing what you lived and felt.
\end{grammarlens}

\subsection*{Your turn}
\begin{itemize}
\raggedright
  \item \emph{Say it:} \emph{virōtam}
  \item \emph{Say it:} \emph{virōtam}, once more
  \item \emph{Say it:} \emph{oru ciṟiya caṇṭaiyāl avarkaḷukkuḷ virōtam vaḷarntatu.} — Because of a small quarrel, hostility grew between them.
  \item \emph{Recall:} say the Tamil for to nod off, then the Tamil for enmity, then say \emph{virōtam} again
\end{itemize}

\begin{tcolorbox}[breakable,colback=teal!4,colframe=teal!35!black,title={Before you move on}]
What does \ta{விரோதம்} mean? (\textquotedblleft{}Hostility; ill will\textquotedblright{}.) Say it once more.
\end{tcolorbox}
\section[\texorpdfstring{\ta{வஞ்சம்} (vañcam)}{வஞ்சம் (vañcam)}]{\ta{வஞ்சம்} (vañcam) — a grudge; spite}

\label{lesson:TA-C469-vancam}



\begin{quote}
Before the new one: say the Tamil for enmity, then the Tamil for hostility.
\end{quote}

\subsection*{You'll want to know: \ta{வஞ்சம்}}
\textbf{\ta{வஞ்சம்}} — \emph{vañcam} — \textquotedblleft{}a grudge; spite\textquotedblright{}.

\begin{quote}
\ta{அவன்} \ta{மனதில்} \ta{பழைய} \ta{வஞ்சம்} \ta{இருந்தது}.
\end{quote}

\begin{quote}
\emph{avaṉ maṉatil paḻaiya vañcam iruntatu.}
\end{quote}

\begin{quote}
He bore an old grudge.
\end{quote}

\ta{வஞ்சம்} \ta{தீர்}, to settle a score.

\begin{cousinweb}
\textbf{\ta{வஞ்சம்}} (\emph{vañcam}) is usually traced to Sanskrit \emph{vañcana}, \textquotedblleft{}cheating, deception\textquotedblright{} (Monier-Williams); in Tamil it is a deceit-filled grudge.
\end{cousinweb}

\begin{grammarlens}[title={What you've built}]
One more word for describing what you lived and felt.
\end{grammarlens}

\subsection*{Your turn}
\begin{itemize}
\raggedright
  \item \emph{Say it:} \emph{vañcam}
  \item \emph{Say it:} \emph{vañcam}, once more
  \item \emph{Say it:} \emph{avaṉ maṉatil paḻaiya vañcam iruntatu.} — He bore an old grudge.
  \item \emph{Recall:} say the Tamil for enmity, then the Tamil for hostility, then say \emph{vañcam} again
\end{itemize}

\begin{tcolorbox}[breakable,colback=teal!4,colframe=teal!35!black,title={Before you move on}]
What does \ta{வஞ்சம்} mean? (\textquotedblleft{}A grudge; spite\textquotedblright{}.) Say it once more.
\end{tcolorbox}
\section[\texorpdfstring{\ta{எதிரி} (etiri)}{எதிரி (etiri)}]{\ta{எதிரி} (etiri) — an enemy; an opponent}

\label{lesson:TA-C469-etiri}



\begin{quote}
Before the new one: say the Tamil for hostility, then the Tamil for a grudge.
\end{quote}

\subsection*{You'll want to know: \ta{எதிரி}}
\textbf{\ta{எதிரி}} — \emph{etiri} — \textquotedblleft{}an enemy; an opponent\textquotedblright{}.

\begin{quote}
\ta{அவன்} \ta{என்} \ta{நண்பன்}, \ta{எதிரி} \ta{இல்லை}.
\end{quote}

\begin{quote}
\emph{avaṉ eṉ naṇpaṉ, etiri illai.}
\end{quote}

\begin{quote}
He is my friend, not my enemy.
\end{quote}

Plural: \ta{எதிரிகள்}, \emph{etirikaḷ}.

From \ta{எதிர்} \textquotedblleft{}opposite\textquotedblright{}: \ta{எதிரே}, opposite, shares it.

\begin{cousinweb}
\textbf{\ta{எதிரி}} (\emph{etiri}) is native Dravidian, built on \ta{எதிர்} \emph{etir}, \textquotedblleft{}that which is opposite, in front\textquotedblright{} (DEDR 795).
\end{cousinweb}

\begin{grammarlens}[title={What you've built}]
One more word for describing what you lived and felt.
\end{grammarlens}

\subsection*{Your turn}
\begin{itemize}
\raggedright
  \item \emph{Say it:} \emph{etiri}
  \item \emph{Say it:} \emph{etiri}, once more
  \item \emph{Say it:} \emph{avaṉ eṉ naṇpaṉ, etiri illai.} — He is my friend, not my enemy.
  \item \emph{Recall:} say the Tamil for hostility, then the Tamil for a grudge, then say \emph{etiri} again
\end{itemize}

\begin{tcolorbox}[breakable,colback=teal!4,colframe=teal!35!black,title={Before you move on}]
What does \ta{எதிரி} mean? (\textquotedblleft{}An enemy; an opponent\textquotedblright{}.) Say it once more.
\end{tcolorbox}
\section[\texorpdfstring{\ta{பழிவாங்கு} (paḻivāṅku)}{பழிவாங்கு (paḻivāṅku)}]{\ta{பழிவாங்கு} (paḻivāṅku) — to take revenge}

\label{lesson:TA-C469-palivanku}



\begin{quote}
Before the new one: say the Tamil for a grudge, then the Tamil for an enemy.
\end{quote}

\subsection*{You'll want to know: \ta{பழிவாங்கு}}
\textbf{\ta{பழிவாங்கு}} — \emph{paḻivāṅku} — \textquotedblleft{}to take revenge\textquotedblright{}.

\begin{quote}
\ta{அந்த} \ta{அவமானத்துக்காக} \ta{அவன்} \ta{பழிவாங்க} \ta{நினைத்தான்}.
\end{quote}

\begin{quote}
\emph{anta avamāṉattukkāka avaṉ paḻivāṅka niṉaittāṉ.}
\end{quote}

\begin{quote}
He thought of taking revenge for that humiliation.
\end{quote}

Past (I …): \ta{பழிவாங்கினேன்}, \emph{paḻivāṅkiṉēṉ}.

\ta{பழி} \textquotedblleft{}blame, vengeance\textquotedblright{} + \ta{வாங்கு} \textquotedblleft{}to take\textquotedblright{}.

\begin{cousinweb}
\textbf{\ta{பழிவாங்கு}} (\emph{paḻivāṅku}) is a native compound of \ta{பழி} \emph{paḻi}, \textquotedblleft{}blame, reproach\textquotedblright{} (DEDR 4002), and \ta{வாங்கு}, \textquotedblleft{}to take\textquotedblright{}.
\end{cousinweb}

\begin{grammarlens}[title={What you've built}]
One more word for describing what you lived and felt.
\end{grammarlens}

\subsection*{Your turn}
\begin{itemize}
\raggedright
  \item \emph{Say it:} \emph{paḻivāṅku}
  \item \emph{Say it:} \emph{paḻivāṅku}, once more
  \item \emph{Say it:} \emph{anta avamāṉattukkāka avaṉ paḻivāṅka niṉaittāṉ.} — He thought of taking revenge for that humiliation.
  \item \emph{Recall:} say the Tamil for a grudge, then the Tamil for an enemy, then say \emph{paḻivāṅku} again
\end{itemize}

\begin{tcolorbox}[breakable,colback=teal!4,colframe=teal!35!black,title={Before you move on}]
What does \ta{பழிவாங்கு} mean? (\textquotedblleft{}To take revenge\textquotedblright{}.) Say it once more.
\end{tcolorbox}
